\documentclass[10pt,a4paper]{amsart}
\usepackage[margin=19mm]{geometry}
\usepackage{amsmath,amssymb,mathtools}
\usepackage[scr=rsfs,cal=euler]{mathalfa}
\usepackage{xfrac}
\usepackage[unicode,hidelinks,bookmarks=false]{hyperref}
\newcommand{\ie}{i.e.\ }
\newcommand{\cf}{cf.\ }
\newcommand{\resp}{respectively\ }
\newcommand{\Subsection}{\subsection}
\providecommand{\nobreakdash}{\nobreak\hbox{-}}
\setlength{\emergencystretch}{2em}
\multlinegap0pt
\usepackage{bm}
\usepackage{bbm}
\usepackage{dsfont}
\usepackage{xspace}
\usepackage{cancel}
\usepackage{mathtools}
\usepackage{amsthm}
\makeatletter
\@ifundefined{theorem}{\newtheorem{theorem}{Theorem}}{}
\@ifundefined{lemma}{\newtheorem{lemma}[theorem]{Lemma}}{}
\makeatother
\makeatletter
\expandafter\ifx\csname proof\endcsname\relax
\newenvironment{proof}{\noindent {\bf Proof.}}{\rule{2mm}{2mm}\par\medskip}
\fi
\expandafter\ifx\csname proofof\endcsname\relax
\newenvironment{proofof}{\noindent {\bf Proof of the Theorem 6.1.}}{\rule{2mm}{2mm}\par\medskip}
\fi
\makeatother
\title[A spectral Erd\H{o}s-Faudree-Rousseau theorem]{A spectral Erd\H{o}s-Faudree-Rousseau theorem}
\author{Yongtao Li, Lihua Feng, Yuejian Peng}
\date{Source: arXiv:2406.13176v2 (CC BY 4.0)}
\begin{document}
\maketitle

\noindent\emph{Source and licence.} A spectral Erd\H{o}s-Faudree-Rousseau theorem. Version arXiv:2406.13176v2 (CC BY 4.0), distributed under the Creative Commons Attribution 4.0 International licence, as recorded in the arXiv licence metadata for this exact version and shown on the abstract page https://arxiv.org/abs/2406.13176v2. Full rights notice: licenses/arxiv-2406.13176-CC-BY-4.0.txt.\par\smallskip

\noindent\textit{Editorial scope.} Counted unit: the Lemma whose source label is \texttt{lem-W-sub-L} in the article (source0.tex lines 1404--1406, source label \texttt{lem-W-sub-L}), delivered together with the article's own complete original proof of it (source0.tex lines 1408--1456, one \texttt{proof} environment of 49 lines). No mathematics is written, repaired or re-derived: statement and proof are the article's own text line for line. Required context retained uncounted: Lupper (lines 1344--1347); set (lines 1396--1401). Every definition, display and label of those lines is retained unchanged. Omitted from this package and not redistributed: 1--1343, 1348--1395, 1402--1403, 1407--1407, 1457--2655. Nothing inside a retained excerpt is omitted or altered: every retained body is delivered exactly as the article's lines are written, so no omission receipt of any kind accompanies this package. Citations: the retained excerpts carry no citation command at all, so no bibliography entry of the article is transcribed and no reference is redistributed. Reference adaptation: none was needed and none was made; every reference inside the delivered unit resolves inside it, so no unresolved reference or citation is produced. Printed slips kept literal: no printed word, symbol, hypothesis or number of the counted statement or of its proof was altered. Layout and class: the article's own class is \texttt{article} with packages including amsmath, amsfonts, color, float, pifont, amssymb, graphicx, bm, enumerate, amsthm, amscd, xy, mdframed, hyperref; the wrapper is the lane's pinned \texttt{amsart} editorial wrapper at 10pt on A4, which restates the theorem-like environments the retained text needs and adds the article's own macro definitions that the retained text needs (none), with extra wrapper packages (bm, bbm, dsfont, xspace, cancel, mathtools). The delivered numbering is the unit's own. Assets: no retained span contains a figure environment, a graphics-inclusion command, a PDF-page inclusion, a table or a TikZ picture; no image file is copied into this package, the package declares no asset dependency and no image file is redistributed with it. Licence: version arXiv:2406.13176v2 of this article is distributed under the Creative Commons Attribution 4.0 International licence, as recorded in the arXiv licence metadata for this exact version and shown on the abstract page https://arxiv.org/abs/2406.13176v2. A full rights notice is delivered at licenses/arxiv-2406.13176-CC-BY-4.0.txt. Characters: every retained excerpt is pure ASCII, so the delivered engine, XeLaTeX with its default Latin Modern text font, reports no missing character.
\begin{lemma}\label{Lupper}  
We have 
$|L| < 10.$ 
\end{lemma}

\begin{lemma}\label{set}
Let $A_1, A_2,\dots, A_k$ be $k$  finite sets. Then
\begin{equation*}\label{set1}
\left|\bigcap_{i=1}^k A_i \right| \ge \sum_{i=1}^k |A_i|-(k-1)\left|\bigcup_{i=1}^k A_i\right|.
\end{equation*}
\end{lemma}

  \begin{lemma} \label{lem-W-sub-L}
  We have $W \subseteq L$ and $|W| \le |L| < 10$. 
  \end{lemma}

  \begin{proof}
We shall prove that if $u\notin L$, then $u\notin W$. 
We denote $L_1=L\cap S$ and $L_2=L\cap T$. Without loss of generality, we may assume that $u\in S$ and 
   $u\notin L_1.$ Since  $S$ and $T$ form a maximum cut in $G$, 
   we claim that $d_T(u)\ge \frac{1}{2}d(u)$. 
   Otherwise, if $d_T(u)< \frac{1}{2}d(u)$, then 
   by $d(u)=d_S(u) + d_T(u)$, we have 
   $d_S(u) > d_T(u)$. Moving the vertex $u$ from 
   $S$ to $T$ yields a new vertex bipartition with more edges, 
   which contradicts with the maximality of $G[S,T]$. 
   So we must have $d_T(u)\ge \frac{1}{2}d(u)$. 
  On the other hand, 
  we have $d(u)> \left(\frac{1}{2}-\frac{1}{200} \right)n$ 
  since $u\not\in L$. Then 
   $$d_T(u)\ge \frac{1}{2}d(u)> 
   \left(\frac{1}{4}-\frac{1}{400} \right)n.$$
  Recall that 
  $|L|< 10$ and $|W| < 1680/ \sqrt{n}$, 
we have 
$  |S\setminus (W\cup L)| \approx \frac{n}{2}$. 
  We claim that $u$ has at most $7$ neighbors in 
  $S\setminus (W\cup L)$. 
  Indeed, suppose on the contrary that 
   $u$ is adjacent to $8$ vertices  $u_1, u_2,\ldots, u_8$ in $S\setminus (W\cup L)$.  Since $u_i\not\in L$,  we have $d(u_i)> \left(\frac{1}{2}-\frac{1}{200}\right )n$. Similarly, 
   we have $d_S(u_i)< \frac{n}{140}$ as $u_i\notin W$. So
  \[  d_T(u_i)=d(u_i)-d_S(u_i)> \left(\frac{1}{2}-\frac{1}{200} - \frac{1}{140} \right)n. \] 
   By Lemma~\ref{set}, we have
   \begin{eqnarray*}
&& \left|N_T(u)\cap N_T(u_1) \cap \cdots 
\cap N_T(u_8) \right|\\[2mm]
 &\ge&  |N_T(u)|+ |N_T(u_1)| + \cdots + |N_T(u_8)|  -8|T| \\
& > &\left (\frac{1}{4}-\frac{1}{400}\right)n+\left(\frac{1}{2}-\frac{1}{200} -\frac{1}{140} \right)n\cdot 8 - 8 \left( \frac{n}{2} + 3n^{1/4}\right)\\
 &>&\frac{n}{9}, 
\end{eqnarray*}
 where the last inequality holds for $n\ge 5191$.  
Let $B$ be the set of common neighbors of $u,u_1,\ldots ,u_8$ in $T$. 
Then $|B|>n/9$. 
Observe that for each vertex $v\in B$, the $vuu_i$ forms a triangle 
for each $1\le i\le 8$, so $vu,vu_i (1\le i \le 8)$ are triangular edges. That is to say, each vertex of $B$ 
is incident to at least $9$ triangular edges. 
This leads to more than $9|B| +8 > n +8$ triangular edges, a contradiction.
Therefore $u$ is adjacent to at most $7$ vertices in  $S\setminus (W\cup L)$. 
Recall that $|L|\le 9$ by Lemma \ref{Lupper}.  
Hence, for $n\ge 5432$, we have 
\[  d_S(u) \le |W|+|L|+ 7 
<  \frac{1680}{\sqrt{n}} + 16 < \frac{n}{140}. \] 
By definition, we get $u\notin W$.
This completes the proof. 
 \end{proof}

\end{document}
